\begin{abstract}
We continue the exact-identity program in the ProveIt manuscript and its
incoming research reports at revision
\texttt{5a790187c8e186e41e2b990b4941cb7a1a3c7b6b}.
Four related questions are treated with complete proofs.
First, the independent-order harmonic normalization is identified with
the established interpolant of Ihara, Nakamura, and Yamamoto. An explicit
prefix-residue proof gives its two-endpoint version, joint entire
continuation in all orders, composition and shift laws, Gamma generating
functions, Stieltjes order jets, and normalized primitives.
Second, every admissible third derivative of
$T(at,bt,ct)$ at zero is reduced to ordinary zeta derivatives and two
specified scalar coordinates: the diagonal derivative $\Omega$ and an
absolutely convergent logarithmic Gamma series $\kappa$. Exact conic and
cyclic combinations eliminate the latter; the full boundary germ is
also evaluated by a convergent Gamma series.
Third, arbitrary weighted colored sums with any mixture of positive
and nonpositive integral inner orders admit a terminating finite
reduction to nested polylogarithms with one free complex outer order.
The construction handles coincident colors, preserves the cumulative
alphabet, controls depth, and gives all Laurent constants at
nonpositive outer orders by finite local coefficient extraction.
Finally, a sampling uniqueness theorem proves the exact finite
single-centre reduction criterion for any number of shifted powers,
including identities on integer Fourier modes.
We distinguish repository advances from known theorems and from
unevaluated coordinates. The Gaussian $S_6$ and revised $S_8$ candidates
remain conjectural. Exact algebra, independent numerical diagnostics,
an executable mixed-order compiler, integration notes, and specific
further research questions accompany the article.
\end{abstract}

\noindent\textbf{Keywords.}
Multiple polylogarithms; Mordell--Tornheim zeta functions; independent
spectral orders; truncated multiple zeta interpolation; generalized
Stieltjes constants; Gamma functions; harmonic sums; Laurent constants;
Fourier uniqueness.
